\documentclass[11pt,a4paper]{article}
\usepackage[T1]{fontenc}
\usepackage[utf8]{inputenc}
\usepackage{lmodern,microtype}
\usepackage[margin=26mm,headheight=15pt]{geometry}
\usepackage{amsmath,amssymb,amsthm,mathtools,bm}
\usepackage{booktabs,array,longtable,enumitem}
\usepackage{xcolor}
\definecolor{inkblue}{HTML}{193C5B}
\usepackage[colorlinks=true,linkcolor=inkblue,citecolor=inkblue,urlcolor=inkblue]{hyperref}
\usepackage{fancyhdr}
\pagestyle{fancy}\fancyhf{}
\fancyhead[L]{\small The near-linear boundary of exponential feedback}
\fancyhead[R]{\small Research article}
\fancyfoot[C]{\thepage}
\setlength{\parskip}{4pt}
\setlength{\parindent}{1.2em}
\setlength{\emergencystretch}{2em}
\linespread{1.06}
\numberwithin{equation}{section}
\newtheorem{theorem}{Theorem}[section]
\newtheorem{proposition}[theorem]{Proposition}
\newtheorem{lemma}[theorem]{Lemma}
\newtheorem{corollary}[theorem]{Corollary}
\theoremstyle{definition}
\newtheorem{definition}[theorem]{Definition}
\newtheorem{example}[theorem]{Example}
\theoremstyle{remark}
\newtheorem{remark}[theorem]{Remark}
\newcommand{\Uhat}{\widehat U}
\newcommand{\Bcal}{\mathcal B}
\newcommand{\R}{\mathbb R}
\newcommand{\N}{\mathbb N}
\newcommand{\eps}{\varepsilon}
\newcommand{\coef}[2]{[#1^{#2}]}
\newcommand{\Prob}{\mathbb P}
\newcommand{\ee}{\mathrm e}
\newcommand{\dd}{\,\mathrm d}
\newcommand{\asympeq}{\mathrel{\sim}}
\DeclareMathOperator{\supp}{supp}
\title{\textbf{The Near-Linear Boundary of\newline Exponential Feedback}\\[7pt]
\large A self-consistent Lambert law, coefficient large deviations,\\
and a cascade of Borel-growth resonances}
\author{Research study prepared for Vladimir Reshetnikov's ProveIt program}
\date{29 September 2026}
\begin{document}
\maketitle
\thispagestyle{empty}
\begin{abstract}
Consider the uniquely defined formal solution
$\Uhat(q)=\sum_{j\ge1}c_jq^j\exp(\lambda_j\Uhat(q))$.
The ProveIt feedback article leaves open the sharp growth law at the
near-linear boundary $\lambda_j=jL(j)$, where $L$ is slowly varying and
unbounded: the solution diverges, although it belongs to every positive
Gevrey class. Under explicit smoothness, convexity, and positive-amplitude
hypotheses, we resolve this coefficient-growth problem. If $a=c_1$ and
$b_n$ is determined by
\[
 b_n\ee^{b_n}=L\!\left(\frac{n}{1+b_n}\right),
\]
then $u_n^{1/n}\sim a\ee^{b_n}$. The Lagrange composition length satisfies
a large-deviation principle with rate $r-1-\log r$ after scaling its
excess by $n/(1+b_n)$. Exponentially controlled amplitudes and bounded
linear perturbations of the slopes preserve these laws.

For every $\sigma>0$, the factorial-normalized transform
$\Bcal_\sigma(t)=\sum u_nt^n/(n!)^\sigma$ is entire, and
$\log\Bcal_\sigma(t)\sim\sigma\nu_\sigma(t)$, where
$\nu_\sigma^\sigma/(\ee^{b_{\nu_\sigma}})=at$.
This identifies a positive-direction obstruction to every finite-order
Borel--Laplace summation procedure of the standard Gevrey kind.
For $L(j)=\exp((\log j)^\gamma)$, $0<\gamma<1$, an explicit
power--logarithmic reversion gives all nonvanishing terms in the logarithm
of the growth normalizer. At $\gamma=1-1/m$, a term becomes constant and
produces the factor $\exp(1/m)$ in ordinary Borel growth. A separate moving-
argument example shows why replacing the implicit Lambert law by
$W(L(n))$ can lose a constant or an unbounded correction. Complete proofs,
finite coefficient bounds, reproducible exact checks, and ten further
research directions are included.
\end{abstract}
\vfill
\noindent\textbf{Scope and status.} The main results are proved here in ordinary
mathematics. They answer a stated repository question under the hypotheses
specified below; they are not a claim to settle an unrestricted problem in
transseries theory. No new Lean proof is supplied. Classical inversion and
regular-variation tools are credited, and publication-level priority has
not been established by an exhaustive literature search.
\clearpage
\tableofcontents
\clearpage

\section{The repository question and the contribution}
\label{sec:context}

The relevant ProveIt material now lives under
\texttt{Analysis/Transseries/}. Its main volume develops power--logarithmic
calculus, formal reversion, exact cores, residual transport, and examples
of special-function inversion. The separate feedback article
\cite{repo-feedback} studies the particularly transparent countable kernel
\begin{equation}
 \Uhat(q)=\sum_{j\ge1}q^j\exp(\lambda_j\Uhat(q)).
 \label{eq:unit-model}
\end{equation}
It supplies a finite positive coefficient formula, a sharp convergence
criterion, a Gevrey classification, and logarithmic coefficient
asymptotics for regularly varying slopes of index strictly greater than
one. Its section entitled \emph{The near-linear transition} explicitly
asks for sharp coefficient asymptotics and natural summation scales when
$\lambda_j=j(\log j)^\beta$ or, more generally, $jL(j)$ with unbounded
slowly varying $L$. The difficulty is not formal existence. It is the
change of scale when the positive multiple of $n\log n$ in the
superlinear analysis disappears.

This article treats that question directly. It also answers a restricted
but useful part of the same source's amplitude question. We allow
\begin{equation}
 \Uhat(q)=\sum_{j\ge1}c_jq^j\exp(\lambda_j\Uhat(q)),
 \qquad c_j>0,\quad\lambda_j\ge0.
 \label{eq:model}
\end{equation}
The central object is not $W(L(n))$ but a Lambert equation with a
\emph{moving argument}. This distinction is necessary at the precision
of an equivalent for $u_n^{1/n}$, rather than merely an equivalent for
$\log u_n$.

\subsection{What is established}
The main coefficient theorem gives
\begin{equation}
 \log u_n=n\bigl(\log a+b_n+o(1)\bigr),\qquad a=c_1.
 \label{eq:headline}
\end{equation}
It is accompanied by a complete large-deviation law for the number of
primitive factors in the finite Lagrange expansion. A second theorem
transfers this precise coefficient scale to the growth of all
factorial-normalized entire transforms. A third theorem computes the
resulting Borel-growth transseries for stretched-logarithmic slopes,
including its constant terms at every transition value $1-1/m$.

These are stronger than a determination of the least Gevrey index, but
we distinguish them carefully from a multiplicative equivalent for
$u_n$ itself. Equation~\eqref{eq:headline} does \emph{not} assert
$u_n\sim(a\ee^{b_n})^n$: an $\exp(o(n))$ factor remains unresolved.
Similarly, the large-deviation result is not a central limit theorem,
and a growth-adapted integral is not automatically an acceleration
operator compatible with the original equation.

\subsection{Relation to established tools and novelty}
The finite coefficient identity is classical Lagrange inversion
\cite{gessel}, already specialized to unit amplitudes in the repository.
General inversion with infinitely many colours has a much broader
literature; Jansen, Kuna, and Tsagkarogiannis \cite{jkt} are a relevant
reference, not a source for the particular asymptotic laws proved here.
Slow variation and its use in inversion are standard themes of
regular-variation theory \cite{bgt}. We use a differentiable subclass so
that the needed estimates can be proved directly rather than imported
as a large technical package.

The new work relative to the inspected repository is the borderline
Lambert normalizer, its probability law, its positive-amplitude and
slope-perturbation stability, and its consequences for Borel growth and
moving power--logarithmic terms. The coefficients in the final scalar
reversion are classical Lagrange coefficients; we do not claim their
formula as a new inversion theorem. Nor do we claim an exhaustive global
priority result for the combined conclusions. The exact source snapshot
and claim boundaries are recorded in Appendix~\ref{app:sources}.

\section{Hypotheses, normalizers, and theorem statements}
\label{sec:statements}

\begin{definition}[Admissible near-linear slopes]
\label{def:admissible}
An admissible slope profile is a positive, nondecreasing $C^1$ function
$L:[1,\infty)\to(0,\infty)$ such that
\begin{equation}
 L(x)\longrightarrow\infty,\qquad
 \eps(x):=\frac{xL'(x)}{L(x)}\longrightarrow0,
 \label{eq:elasticity}
\end{equation}
and $\Lambda(x)=xL(x)$ is convex. The base slopes are
$\lambda_j=\Lambda(j)$. More generally, the perturbation theorem allows
nonnegative $\widetilde\lambda_j$ with
\begin{equation}
 |\widetilde\lambda_j-\Lambda(j)|\le Cj
 \quad(j\ge1)
 \label{eq:linear-perturbation}
\end{equation}
for a fixed $C$. Conditions imposed only for large $x$ may be used after
choosing an admissible extension; finite discrepancies are covered by
\eqref{eq:linear-perturbation}.
\end{definition}

Condition~\eqref{eq:elasticity} implies slow variation: integrate
$L'/L$ between $x$ and $rx$ to obtain $L(rx)/L(x)\to1$ for each fixed
$r>0$. It also gives $\log L(x)=o(\log x)$.

\begin{definition}[Exponentially controlled positive amplitudes]
\label{def:amplitudes}
Put $a=c_1>0$ and $d_j=c_j/a^j$. We require a constant $D\ge0$ such that
\begin{equation}
 \ee^{-D(j-1)}\le d_j\le\ee^{D(j-1)}\quad(j\ge1).
 \label{eq:amplitudes}
\end{equation}
Thus $d_1=1$. Polynomial amplitudes, positive constant amplitudes,
positive geometric amplitudes, and finite modifications of such families
are included. Superexponentially small amplitudes and missing actions are
not included.
\end{definition}

For all sufficiently large real $x$, define $b(x)>0$ by
\begin{equation}
 b(x)\ee^{b(x)}=L\!\left(\frac{x}{1+b(x)}\right),\qquad
 f(x)=\ee^{b(x)},\qquad x_*(x)=\frac{x}{1+b(x)}.
 \label{eq:b-definition}
\end{equation}
We abbreviate $b_n=b(n)$ and $x_n=x_*(n)$. Existence, uniqueness, and
regularity are proved in Section~\ref{sec:normalizer}.

\begin{theorem}[Sharp borderline coefficient law]
\label{thm:main}
Under Definitions~\ref{def:admissible} and \ref{def:amplitudes}, including
the perturbation option~\eqref{eq:linear-perturbation}, the unique formal
solution $\Uhat=\sum_{n\ge1}u_nq^n$ has positive coefficients and satisfies
\begin{equation}
 \boxed{\quad u_n^{1/n}\sim a f(n)=a\ee^{b_n}.\quad}
 \label{eq:root-law}
\end{equation}
In particular, its radius of convergence is zero, while for every
$\sigma>0$ and every $A>0$ there is a finite $C_{\sigma,A}$ such that
\begin{equation}
 u_n\le C_{\sigma,A} A^n(n!)^\sigma\quad(n\ge1).
 \label{eq:zero-index}
\end{equation}
\end{theorem}

The arbitrary $A>0$ in \eqref{eq:zero-index} says more than positive
Gevrey membership: the factorial-normalized type is zero at every
positive order. The series still fails to be convergent.

For the probabilistic refinement, set
\begin{equation}
 W_{n,k}=\frac1{k!}
 \sum_{\substack{j_1+\cdots+j_k=n\\j_i\ge1}}
 \left(\prod_{i=1}^k c_{j_i}\right)
 \left(\sum_{i=1}^k\lambda_{j_i}\right)^{k-1}.
 \label{eq:Wnk}
\end{equation}
As usual, a zero exponent is interpreted as one, including $0^0$ here.
Then $u_n=\sum_kW_{n,k}$. Give $K_n$ the distribution
$\Prob(K_n=k)=W_{n,k}/u_n$, and define
\begin{equation}
 R_n=\frac{n-K_n+1}{x_n},\qquad I(r)=r-1-\log r\quad(r>0).
 \label{eq:Rn}
\end{equation}

\begin{theorem}[Length large deviations]
\label{thm:ldp}
The variables $R_n$ obey a large-deviation principle on $(0,\infty)$
with speed $n$ and good rate function $I$. Explicitly, for open $G$ and
closed $F$ in $(0,\infty)$,
\begin{align}
 \liminf_{n\to\infty}\frac1n\log\Prob(R_n\in G)
 &\ge-\inf_{r\in G}I(r),\\
 \limsup_{n\to\infty}\frac1n\log\Prob(R_n\in F)
 &\le-\inf_{r\in F}I(r).
 \label{eq:LDP}
\end{align}
If $(n-k_n+1)/x_n\to r\in(0,\infty)$, then the local statement is
\begin{equation}
 \frac1n\log\frac{W_{n,k_n}}{u_n}\longrightarrow-I(r).
 \label{eq:local-LDP}
\end{equation}
Consequently $R_n\to1$ in probability, exponentially at every fixed
positive distance from one.
\end{theorem}

\begin{theorem}[Factorial-Borel growth]
\label{thm:borel}
For $\sigma>0$, the function
\begin{equation}
 \Bcal_\sigma(t)=\sum_{n\ge1}\frac{u_nt^n}{(n!)^\sigma}
 \label{eq:Borel}
\end{equation}
is entire. Let $\nu_\sigma(t)$ be the eventually unique large positive
solution of
\begin{equation}
 \frac{\nu_\sigma(t)^\sigma}{f(\nu_\sigma(t))}=at.
 \label{eq:nu-definition}
\end{equation}
Then, on the positive real ray,
\begin{equation}
 \boxed{\quad\log\Bcal_\sigma(t)\sim\sigma\nu_\sigma(t).\quad}
 \label{eq:Borel-growth}
\end{equation}
The maximum modulus equals $\Bcal_\sigma(t)$ on the circle $|z|=t$.
Its order is $1/\sigma$ and its type at that order is infinite.
\end{theorem}

The theorem gives an implicit asymptotic equivalent, not merely an order
of growth. It is that extra precision which detects the constant factors
in Section~\ref{sec:resonance}.

\section{Formal existence and finite coefficient inequalities}
\label{sec:finite}

\begin{proposition}[Formal locality and coefficient formula]
\label{prop:formal}
Equation~\eqref{eq:model} has a unique solution in $q\R[[q]]$.
The coefficient $u_n$ depends only on the actions and slopes with index
at most $n$, and it is given by $u_n=\sum_{k=1}^nW_{n,k}$.
\end{proposition}
\begin{proof}
The map $V\mapsto\sum_jc_jq^j\exp(\lambda_jV)$ raises the $q$-adic
valuation of a difference by at least one. Starting at zero therefore
stabilizes every coefficient and gives a unique fixed point. In degree
$n$, only $j\le n$ and finitely many exponential terms enter.

For the formula, introduce a marker $z$ and solve
$V=z\sum_jc_jq^j\exp(\lambda_jV)$. Formal Lagrange inversion gives
\[
 [z^k]V=\frac1k[t^{k-1}]
 \left(\sum_jc_jq^j\exp(\lambda_jt)\right)^k.
\]
The ordered composition $(j_1,\ldots,j_k)$ contributes the coefficient
of $t^{k-1}$ in $\exp((\sum_i\lambda_{j_i})t)$, proving
\eqref{eq:Wnk}. Specializing $z=1$ is coefficientwise legitimate because
$k\le n$ in degree $q^n$. Positivity follows from the $k=1$ term $c_n>0$.
\end{proof}

For example, without any asymptotic hypotheses,
\begin{align*}
 u_1&=c_1,\\
 u_2&=c_2+c_1^2\lambda_1,\\
 u_3&=c_3+c_1c_2(\lambda_1+\lambda_2)
            +\tfrac32c_1^3\lambda_1^2.
\end{align*}
The $k$-marker will be important: a formal coefficient formula does not
by itself say which composition lengths supply its mass.

\begin{lemma}[Convex composition bound]
\label{lem:convex}
If the base sequence $\lambda_j=\Lambda(j)$ is convex, then for
$j_1+\cdots+j_k=n$,
\begin{equation}
 \sum_{i=1}^k\lambda_{j_i}
 \le\lambda_{n-k+1}+(k-1)\lambda_1.
 \label{eq:convex-bound}
\end{equation}
\end{lemma}
\begin{proof}
The forward differences $\lambda_{j+1}-\lambda_j$ are nondecreasing.
If $r\ge s>1$, moving one unit from $s$ to $r$ cannot decrease the sum
$\lambda_r+\lambda_s$. Repeating the move leaves at most one part larger
than one. Its size is $n-k+1$, which proves the inequality.
\end{proof}

Write $m=n-k+1$ and
\begin{equation}
 A_{n,k}=\lambda_m+(k-1)\lambda_1,\qquad
 E_{n,k}=\frac{A_{n,k}^{k-1}}{k!},\qquad
 r_{n,k}=\begin{cases}k,&m>1,\\1,&m=1.\end{cases}
 \label{eq:finite-envelopes}
\end{equation}
For unit amplitudes, one composition with a single part $m$ and the
remaining parts one has slope sum $A_{n,k}$. It has $k$ distinct
placements unless $m=1$.

\begin{proposition}[Finite, computable coefficient bounds]
\label{prop:finite-bounds}
For the unperturbed convex slopes and amplitudes
\eqref{eq:amplitudes},
\begin{equation}
 a^n\ee^{-D(n-k)}r_{n,k}E_{n,k}
 \le W_{n,k}\le
 a^n\ee^{D(n-k)}\binom{n-1}{k-1}E_{n,k}.
 \label{eq:finite-W-bound}
\end{equation}
In particular,
\begin{equation}
 \max_k a^n\ee^{-D(n-k)}r_{n,k}E_{n,k}
 \le u_n\le
 \sum_k a^n\ee^{D(n-k)}\binom{n-1}{k-1}E_{n,k}.
 \label{eq:finite-u-bound}
\end{equation}
\end{proposition}
\begin{proof}
Each composition satisfies
\[
 \prod_i c_{j_i}=a^n\prod_i d_{j_i},\qquad
 \ee^{-D(n-k)}\le\prod_i d_{j_i}\le\ee^{D(n-k)},
\]
because $\sum_i(j_i-1)=n-k$. Use the exceptional-part compositions for
the lower bound, and \eqref{eq:convex-bound} and the exact number
$\binom{n-1}{k-1}$ of compositions for the upper bound.
\end{proof}

These bounds require only $O(n)$ envelope evaluations once the first $n$
slopes are known. They can be evaluated with directed rounding to obtain
numerical intervals. The supplied floating-point tables do not implement
that rounding; the small rational tests check the inequalities exactly.

\section{The self-consistent Lambert normalizer}
\label{sec:normalizer}

\begin{lemma}[Existence and regularity]
\label{lem:b}
For large $x$, \eqref{eq:b-definition} has exactly one solution with
$0<b(x)<x-1$. Moreover,
\begin{gather}
 b(x)\to\infty,\qquad b(x)=o(\log x),\qquad x_*(x)\to\infty,\\
 b(x)\sim\log L(x),\qquad
 \frac{\dd b(x)}{\dd\log x}
 =\frac{\eps(x_*(x))}
 {1+1/b(x)+\eps(x_*(x))/(1+b(x))}\longrightarrow0.
 \label{eq:b-derivative}
\end{gather}
Thus $f(x)=\ee^{b(x)}$ is nondecreasing, unbounded, and slowly varying.
For each $\sigma>0$, $x^\sigma/f(x)$ is eventually strictly increasing,
unbounded, and regularly varying with index $\sigma$.
\end{lemma}
\begin{proof}
As a function of $b$, the left side $b\ee^b$ is strictly increasing and
the right side $L(x/(1+b))$ is nonincreasing. At $b=0$ the former is zero;
at $b=x-1$ it exceeds $L(1)$ for large $x$. This proves existence and
uniqueness. Since $b\le W(L(x))\le\log L(x)$ once $b\ge1$, we have
$b=o(\log x)$, so $x_*(x)\to\infty$. If $b$ remained bounded along a
subsequence, the defining equation would contradict $L(x_*)\to\infty$.

For large $x$, $0\le\eps\le1$ throughout $[x_*,x]$. Integration gives
\[
 0\le\log L(x)-\log L(x_*)\le\log(1+b).
\]
But $\log L(x_*)=b+\log b$, so $\log L(x)=b+O(\log b)$, proving the
claimed equivalent. Implicit differentiation of
$b+\log b=\log L(x/(1+b))$ gives \eqref{eq:b-derivative}.
Integration of that derivative on fixed logarithmic intervals proves
slow variation of $f$. Finally,
\[
 \frac{\dd}{\dd\log x}\log\frac{x^\sigma}{f(x)}
 =\sigma-\frac{\dd b}{\dd\log x}\longrightarrow\sigma>0,
\]
which proves the remaining assertions.
\end{proof}

The same calculation shows that $\nu_\sigma$ in
\eqref{eq:nu-definition} exists uniquely for large arguments. Its regular
variation index is $1/\sigma$: taking ratios in the defining equation
and using the derivative bound gives this directly.

\subsection{Why the argument moves}
Ignoring the fixed $\lambda_1$ contribution for the moment, the exceptional-
part term with $m$ large has logarithm
\[
 (n-m)\left(\log\frac{mL(m)}{n-m}+1\right)+O(\log n).
\]
Write $y=(n-m)/m$. Its logarithm divided by $n$ becomes
\begin{equation}
 \frac{y}{1+y}
 \left(\log L\!\left(\frac{n}{1+y}\right)-\log y+1\right).
 \label{eq:y-envelope}
\end{equation}
If $y=b_n$, the expression equals \emph{exactly} $b_n$.
This cancellation is the reason for the normalization
\eqref{eq:b-definition}. The actual finite-$n$ maximizer need not have
$y=b_n$; the proof must show that optimizing it changes the normalized
logarithm by only $o(1)$. We do that next, including the full composition
count rather than discarding it.

\section{The borderline envelope and the length law}
\label{sec:envelope}

All statements in this section first concern the unperturbed convex
slopes. Put $b=b_n$, $x_*=n/(1+b)$, $\rho=m/n$, and define
\begin{align}
 H_n(\rho)&=(1-\rho)
 \left(\log\frac{\rho L(n\rho)}{1-\rho}+1\right),
 \label{eq:H}\\
 H_n^{(c)}(\rho)&=(1-\rho)
 \left(\log\left(\frac{\rho L(n\rho)}{1-\rho}+c\right)+1\right),
 \label{eq:Hc}
\end{align}
where $c=\lambda_1$. Endpoint values at $\rho=1$ are understood by
limits. At $\rho=m/n$, Stirling's formula, uniformly in $1\le k\le n$,
gives
\begin{equation}
 \frac1n\log E_{n,k}=H_n^{(c)}(\rho)+o(1),\qquad k=n-m+1.
 \label{eq:stirling-envelope}
\end{equation}
The harmless shifts by one contribute at most $O(\log(nL(n))/n)$.
For $k=1$ the exact left side is zero, consistent with the endpoint.
Also, with $h(\rho)=-\rho\log\rho-(1-\rho)\log(1-\rho)$,
\begin{equation}
 \frac1n\log\binom{n-1}{k-1}=h(\rho)+o(1)
 \label{eq:entropy}
\end{equation}
uniformly. Both uniform estimates follow, for example, from the
factorial bounds $\log(r!)=r\log r-r+O(\log(r+2))$.

\begin{lemma}[Compact profile and coercivity]
\label{lem:profile}
Uniformly for $r$ in compact subsets of $(0,\infty)$,
\begin{equation}
 H_n\!\left(\frac{r}{1+b_n}\right)-b_n
 \longrightarrow\log r+1-r.
 \label{eq:profile}
\end{equation}
For fixed $D\ge0$ and $c\ge0$,
\begin{equation}
 \max_{1\le m\le n}
 \left\{H_n^{(c)}(m/n)+h(m/n)+D\frac mn\right\}
 =b_n+o(1).
 \label{eq:sharp-envelope}
\end{equation}
Moreover the expression minus $b_n$ is coercive in
$r=m/x_n$: along any sequence with $r\to0$ or $r\to\infty$, it tends to
$-\infty$.
\end{lemma}
\begin{proof}
For $\rho=r/(1+b)$, let
$\Delta_n(r)=\log L(rx_*)-\log L(x_*)$.
Since $\log L(x_*)=b+\log b$, a direct calculation gives
\begin{equation}
 H_n(\rho)-b
 =(1-\rho)\left(
 \log\frac{b}{1+b-r}+\log r+\Delta_n(r)+1\right)-b\rho.
 \label{eq:profile-identity}
\end{equation}
On a fixed compact $r$-interval, $\Delta_n(r)\to0$ uniformly by
\eqref{eq:elasticity}, $\rho\to0$, and $b\rho\to r$.
This proves \eqref{eq:profile}.

We give the tail argument explicitly, since a bound by $2^n$ throughout
would leave an erroneous constant of order $n$.
First, if $\rho\ge\eta>0$, then
\[
 H_n(\rho)\le(1-\eta)\log L(n)+O(1)
              \le(1-\eta)b+O(\log b).
\]
Here $(1-\rho)\log(\rho/(1-\rho))$ is bounded above on $(0,1)$.
The entropy and $D\rho$ are bounded, so these terms are far below $b$.

The addition of $c$ causes no problem. Put
$X=\rho L(n\rho)/(1-\rho)$. If $X\ge\ee^{b/2}$, then
\[
 0\le H_n^{(c)}-H_n
 \le\log(1+c\ee^{-b/2})=o(1).
\]
If $X<\ee^{b/2}$, then $H_n^{(c)}\le b/2+O(1)$, which is again far
below $b$ even after entropy and amplitude terms. Therefore every
potentially competitive sequence must have $\rho\to0$ and may be
analyzed using $H_n$; then $h(\rho)+D\rho\to0$.

Suppose now that $\rho\to0$ but $r\to\infty$. Both $x_*$ and $rx_*$
are large. Given $\delta>0$, integration of \eqref{eq:elasticity} gives
$0\le\Delta_n(r)\le\delta\log r$ eventually. In
\eqref{eq:profile-identity},
$\log(b/(1+b-r))=o(1)$, while $b\rho=r(1+o(1))$.
Thus its right side is bounded above by
\[
 (1+\delta)\log r+1-r(1+o(1))+o(1),
\]
which tends to $-\infty$.
If $r\to0$ and $m\to\infty$, monotonicity gives $\Delta_n(r)\le0$;
the $\log r$ term then forces the same conclusion. Bounded $m$ is
even simpler: $H_n^{(c)}$ stays bounded above, or tends downward, while
$b\to\infty$.

These arguments prove coercivity. Any sequence of maximizers is therefore
confined to a compact $r$-interval after discarding a vanishing error.
On such an interval, the limiting profile is $\log r+1-r$, whose unique
maximum is zero at $r=1$. Taking $m$ nearest to $x_n$ gives the matching
lower bound. This proves \eqref{eq:sharp-envelope}.
\end{proof}

\begin{proof}[Proof of Theorems~\ref{thm:main} and \ref{thm:ldp}
for the base slopes]
Use \eqref{eq:finite-u-bound}. The upper logarithm differs from the
maximum of its $n$ summands by at most $\log n=o(n)$.
Equations~\eqref{eq:stirling-envelope}--\eqref{eq:sharp-envelope} show
that this maximum divided by $n$ is $\log a+b_n+o(1)$.
For the lower bound choose $m$ nearest to $x_n$. Then
$D(n-k)/n\to0$, and the compact profile at $r=1$ gives the same value.
This proves \eqref{eq:root-law}.

For $m/x_n\to r\in(0,\infty)$, the upper and lower estimates for
$W_{n,k}$ differ after division by $n$ by at most
$2D(m-1)/n+h(m/n)+o(1)=o(1)$. Hence, uniformly on compact $r$-intervals,
\begin{equation}
 \frac1n\log W_{n,k}-\log a-b_n\longrightarrow-I(r).
 \label{eq:local-W}
\end{equation}
Subtracting the coefficient law proves \eqref{eq:local-LDP}.

For an open set, choose one lattice point $m/x_n$ approaching any of its
points; $x_n\to\infty$ makes this possible. This proves the lower
large-deviation bound. On a compact set, the largest of at most $n$
weights gives the upper bound, since $\log n/n\to0$.
Lemma~\ref{lem:profile}'s coercivity implies exponential tightness: for
every $M>0$ there are $0<\delta<R<\infty$ such that the total probability
outside $[\delta,R]$ is at most $\exp(-Mn+o(n))$.
This extends the upper bound to every closed set and completes the LDP.
Since $I$ vanishes only at one, exponential concentration follows.

Finally $b_n\to\infty$ proves radius zero, and
$b_n=o(\log n)$ combined with Stirling's formula proves
\eqref{eq:zero-index}, including arbitrary $A>0$.
\end{proof}

\begin{remark}[What length concentration does not say]
The variable $n-K_n$ is the total excess $\sum_i(j_i-1)$, not the largest
part. A composition with one exceptional action is enough to prove the
lower envelope, but the LDP does not by itself prove that a typical
weighted composition has a unique giant action. Near index one, the
penalty for splitting a large action can be smaller than order $n$.
Identifying that smaller scale is a separate problem.
\end{remark}

\section{Robustness under linear slope errors}
\label{sec:robustness}

\begin{proposition}[Stability at root and large-deviation precision]
\label{prop:robustness}
The coefficient law and the length LDP remain valid when the slopes are
replaced by any nonnegative sequence satisfying
\eqref{eq:linear-perturbation}. The same normalizer $b_n$ is used.
\end{proposition}
\begin{proof}
For a composition of $n$, let $S=\sum_i\Lambda(j_i)$ and
$\widetilde S=\sum_i\widetilde\lambda_{j_i}$. Then
$|\widetilde S-S|\le Cn$, independently of its length.
Split the finite sum into high terms with $S\ge n\ee^{b_n/2}$ and low
terms with $S<n\ee^{b_n/2}$.

On the high set,
\[
 \frac{\widetilde S}{S}=1+O(\ee^{-b_n/2}),
 \qquad
 \left|\log\frac{\widetilde S^{k-1}}{S^{k-1}}\right|
 \le O(n\ee^{-b_n/2})=o(n)
\]
uniformly. The amplitudes are unchanged.
On the low set, both slope sums are at most
$n(\ee^{b_n/2}+C)$. Because this bound divided by $n$ tends to infinity,
$A^{k-1}/k!$ with $A=n(\ee^{b_n/2}+C)$ increases for $1\le k\le n$,
up to an immaterial first factor. Its maximum has logarithm
$n(b_n/2+1+o(1))$. There are at most $2^{n-1}$ compositions altogether,
and their amplitude product is at most $a^n\ee^{Dn}$.
The total low contribution therefore has logarithm at most
\[
 n\bigl(\log a+b_n/2+D+1+\log2+o(1)\bigr),
\]
which is negligible relative to $n(\log a+b_n)$.

Thus the two total coefficient logarithms differ by $o(n)$.
At every fixed scaled length $r>0$, the exceptional-part lower bound
has $S/n$ of order $\ee^{b_n}$, so it is on the high set. The local LDP
and its compact upper bounds therefore transfer with error $o(n)$.
The same high/low split preserves exponential tightness, proving the
full LDP for the perturbed slopes.
\end{proof}

This permits convenient exact models and harmless changes at small
indices. It does not license a replacement of $L(x_n)$ by $L(n)$:
that would change each slope at large indices by more than a fixed
linear error in some of our examples.

\section{From coefficient growth to entire Borel growth}
\label{sec:borel-proof}

We prove a general maximum-term lemma, making explicit the input needed
from the coefficient theorem.

\begin{lemma}[Slowly varying coefficient transfer]
\label{lem:entire-transfer}
Let $a>0$, let $f$ be positive, nondecreasing, unbounded, and $C^1$ for
large arguments, and assume
$\dd\log f(x)/\dd\log x\to0$. If $v_n>0$ satisfy
$v_n^{1/n}\sim af(n)$, then for $\sigma>0$,
\[
 \log\left(\sum_{n\ge1}\frac{v_nt^n}{(n!)^\sigma}\right)
 \sim\sigma N(t),\qquad N(t)^\sigma/f(N(t))=at.
\]
\end{lemma}
\begin{proof}
Since $f(n)=n^{o(1)}$, the series is entire. Write
$\log v_n=n(\log a+\log f(n)+\eta_n)$, where $\eta_n\to0$, and
put $N=N(t)$. Stirling's formula gives
\begin{equation}
 \log\frac{v_nt^n}{(n!)^\sigma}
 =n\left(\sigma\log\frac Nn+
          \log\frac{f(n)}{f(N)}+\sigma+\eta_n\right)+O(\log(n+2)).
 \label{eq:entire-term}
\end{equation}
For $n/N\to r$ in a compact subinterval of $(0,\infty)$, divide by $N$:
the right side tends uniformly to $\sigma r(1-\log r)$.
Its unique maximum is $\sigma$ at $r=1$. The nearest integer to $N$
therefore supplies the lower bound.

For completeness, control the tails. For $n\le\delta N$, beyond a fixed
initial segment, monotonicity gives $f(n)\le f(N)$, and $|\eta_n|$ is
small. The upper logarithm divided by $N$ is bounded by
$\delta(\sigma\log(1/\delta)+\sigma+o(1))$ when $\delta$ is small.
The fixed initial segment has logarithm $O(\log t)=o(N)$.

For $n\ge RN$, the derivative hypothesis gives, uniformly,
$f(n)/f(N)\le(n/N)^{\sigma/2}$ once $N$ is large. Formula
\eqref{eq:entire-term} is then at most
\[
 n\left(-\frac\sigma2\log(n/N)+\sigma+o(1)\right).
\]
For sufficiently large fixed $R$, these terms have a summable
geometric upper bound, and their total is exponentially below the
central maximum. On $[\delta N,RN]$ there are $O(N)$ terms, so summing
adds only $O(\log N)=o(N)$ to the maximum logarithm. Letting the compact
interval expand proves the result.
\end{proof}

\begin{proof}[Proof of Theorem~\ref{thm:borel}]
Apply Lemma~\ref{lem:entire-transfer} with the coefficient theorem and
Lemma~\ref{lem:b}. Nonnegative coefficients give
$|\Bcal_\sigma(z)|\le\Bcal_\sigma(|z|)$, with equality on the positive
ray. The regular variation of $\nu_\sigma$ has index $1/\sigma$, and
\[
 \frac{\nu_\sigma(t)}{t^{1/\sigma}}
   =\bigl(a f(\nu_\sigma(t))\bigr)^{1/\sigma}\longrightarrow\infty.
\]
These facts prove the order and infinite-type assertions.
\end{proof}

\subsection{Standard Gamma normalization and summability}
For clarity, the standard moment normalization associated with Gevrey
order $\sigma$ is
\begin{equation}
 \mathfrak B_\sigma(t)=
 \sum_{n\ge1}\frac{u_nt^n}{\Gamma(1+\sigma n)}.
 \label{eq:gamma-borel}
\end{equation}
Our $\Bcal_1$ is an unshifted Borel transform; differentiating it gives
the common shifted convention. The obstruction below is independent of
this harmless convention.

\begin{corollary}[No positive-direction standard summation of finite order]
\label{cor:no-laplace}
For every $\sigma>0$, $\mathfrak B_\sigma$ is entire and
\begin{equation}
 \log\mathfrak B_\sigma(t)
 \sim\sigma\nu_\sigma(t/\sigma^\sigma),\qquad t\to+\infty.
 \label{eq:gamma-growth}
\end{equation}
In particular, for every $A>0$,
\begin{equation}
 \int_1^\infty t^{1/\sigma-1}
 \exp(-A t^{1/\sigma})\mathfrak B_\sigma(t)\dd t=+\infty.
 \label{eq:laplace-diverges}
\end{equation}
Thus the formal solution is not summable in the positive direction by
any standard finite-order Gevrey Borel--Laplace procedure.
\end{corollary}
\begin{proof}
Stirling's formula gives
\[
 \log\Gamma(1+\sigma n)-\sigma\log(n!)
 =\sigma n\log\sigma+O(\log(n+2)).
\]
Consequently the maximum-term proof applies with $t$ replaced by
$t/\sigma^\sigma$; the polynomial factors do not affect the leading
logarithm. The quotient of \eqref{eq:gamma-growth} by $t^{1/\sigma}$
tends to infinity. The integrand in \eqref{eq:laplace-diverges}
therefore eventually exceeds an exponentially increasing function,
which proves divergence.
\end{proof}

Borel--Laplace summability requires both continuation and growth control;
see Sauzin \cite{sauzin}. The transforms here already continue to the
entire plane. Their obstruction is growth on the chosen ray, not a finite
singularity. Gevrey implicit-function theorems such as Nikolaev's
\cite{nikolaev} have analytic hypotheses which cannot be replaced by
coefficient membership alone. Nothing in Corollary~\ref{cor:no-laplace}
excludes useful sums in other directions.

\begin{proposition}[A critical positive damping scale]
\label{prop:damping}
For the ordinary Borel transform, put $\nu=\nu_1$. Then
\[
 \int_1^\infty\exp(-A\nu(t))\Bcal_1(t)\dd t
\]
converges for $A>1$ and diverges for $0<A<1$.
\end{proposition}
\begin{proof}
Use $\log\Bcal_1(t)=(1+o(1))\nu(t)$ and
$\nu(t)/t=af(\nu(t))\to\infty$.
\end{proof}
This is a growth statement only. It does not define an acceleration
operator respecting products, differentiation, or the feedback equation.
The borderline $A=1$ also requires information not provided by the
leading equivalent.

\subsection{A literal positive fixed point is impossible}
Under our amplitude condition and $L(j)\to\infty$, for $q>0$ and $u>0$
\[
 \log(c_jq^j\ee^{\lambda_ju})
 \ge j(\log a+\log q-D+uL(j))+O(1)\longrightarrow+\infty
\]
for base slopes, and the same is true after a bounded linear slope
perturbation. Thus the summands do not tend to zero. No proposed positive-
ray regularization can be described as a solution of the literal
convergent positive sum without changing its analytic interpretation.

\section{Explicit profiles and the moving-argument correction}
\label{sec:examples}

The implicit normalizer is already an asymptotic equivalent. To extract
its visible power--logarithmic terms, write
\begin{equation}
 s=\log x,\qquad h(s)=\log L(\ee^s).
 \label{eq:h-profile}
\end{equation}
The notation $h$ here is unrelated to the binary entropy used in
Section~\ref{sec:envelope}.

\begin{proposition}[A moving-argument expansion]
\label{prop:moving}
Suppose the admissible profile has $h\in C^2$ eventually and
\begin{equation}
 (\log h(s))^2
 \sup_{|v-s|\le2\log h(s)}|h''(v)|\longrightarrow0.
 \label{eq:h-second}
\end{equation}
Then
\begin{equation}
 b(\ee^s)=h(s)-\log h(s)-h'(s)\log h(s)+o(1).
 \label{eq:moving-expansion}
\end{equation}
In particular,
\begin{equation}
 b(\ee^s)-W(L(\ee^s))=-h'(s)\log h(s)+o(1).
 \label{eq:W-difference}
\end{equation}
\end{proposition}
\begin{proof}
Lemma~\ref{lem:b} gives $b\sim h(s)$, so
$\log b=\log h(s)+o(1)$ and
$\log(1+b)=\log h(s)+o(1)$. Also $h'(s)=\eps(\ee^s)\to0$.
In the equation
\[
 b+\log b=h\bigl(s-\log(1+b)\bigr),
\]
we may therefore replace the right-hand argument by $s-\log h(s)$ at
an additive cost $o(1)$. Taylor's theorem and \eqref{eq:h-second} give
\[
 h(s-\log h(s))=h(s)-h'(s)\log h(s)+o(1),
\]
which proves \eqref{eq:moving-expansion}. Finally the positive solution
of $w+\log w=h$ satisfies $w=h-\log h+o(1)$: substituting the proposed
value leaves an $o(1)$ residual, and the derivative $1+1/w$ is bounded
below by one. This proves \eqref{eq:W-difference}; the conventional
Lambert notation agrees with DLMF~\S4.13 \cite{dlmf}.
\end{proof}

\subsection{Logarithmic slope amplification}
Let $\beta>0$ and suppose $L(x)=(\log x)^\beta$ for large $x$,
with an admissible extension at small $x$. Then $h(s)=\beta\log s$,
and the derivative correction in \eqref{eq:moving-expansion} tends to
zero. Consequently
\begin{align}
 b_n&=\beta\log\log n-\log\log\log n-\log\beta+o(1),\\
 u_n^{1/n}&\sim
 a\frac{(\log n)^\beta}{\beta\log\log n}.
 \label{eq:log-power-root}
\end{align}
This resolves the explicit $j(\log j)^\beta$ family in the repository's
near-linear question at coefficient-root precision.

For the ordinary Borel transform,
\begin{equation}
 \log\Bcal_1(t)\sim
 a t\frac{(\log t)^\beta}{\beta\log\log t}.
 \label{eq:log-power-borel}
\end{equation}
Indeed, $f(x)\sim(\log x)^\beta/(\beta\log\log x)$ satisfies
$f(atf(t))/f(t)\to1$, so the inverse relation
$\nu_1=atf(\nu_1)$ gives $\nu_1\sim atf(t)$. This last conclusion can
also be proved directly by substituting $(1\pm\eta)atf(t)$ and using
monotonicity, then letting $\eta\downarrow0$.

More generally, putting $R=\sigma^{-1}\log(at)$ yields
\begin{equation}
 \log\Bcal_\sigma(t)\sim
 \sigma(at)^{1/\sigma}
 \left(\frac{R^\beta}{\beta\log R}\right)^{1/\sigma}.
 \label{eq:log-power-general-borel}
\end{equation}
All constants, including the $\sigma$ inside $R$, matter in this formula.

\subsection{An exact rational example}
The slopes
\begin{equation}
 \lambda_j=jH_j,\qquad H_j=\sum_{r=1}^j\frac1r,
 \label{eq:harmonic-slopes}
\end{equation}
are rational and convex: their successive difference is $H_j+1$,
which increases. They differ from $j\log j$ by $O(j)$, so
Proposition~\ref{prop:robustness} applies. With unit amplitudes,
\begin{equation}
 u_n^{1/n}\sim\frac{\log n}{\log\log n},\qquad
 \log\Bcal_1(t)\sim\frac{t\log t}{\log\log t}.
 \label{eq:harmonic-results}
\end{equation}
For a concrete smooth admissible interpolation used in the numerical
envelopes, set $L(x)=\psi(x+1)+\gamma_{\!E}$. The identity
\[
 L(x)=\sum_{k\ge1}\left(\frac1k-\frac1{x+k}\right)
\]
implies $L'>0$ and
\[
 (xL(x))''=2\sum_{k\ge1}\frac{k}{(x+k)^3}>0.
\]
Also $L(x)\sim\log x$ and $xL'(x)/L(x)\to0$, directly by integral
comparison. Thus the interpolation is admissible without an appeal to
any unproved numerical property of the digamma function.

\subsection{Stretched-logarithmic amplification}
For $0<\gamma<1$, let $L(x)=\exp((\log x)^\gamma)$ eventually.
Here $h(s)=s^\gamma$, and Proposition~\ref{prop:moving} gives
\begin{equation}
 b_n=(\log n)^\gamma-\gamma\log\log n+o(1),\qquad
 u_n^{1/n}\sim a\frac{\exp((\log n)^\gamma)}{(\log n)^\gamma}.
 \label{eq:stretch-root}
\end{equation}
Although the moving argument has only an $o(1)$ effect on this particular
$b_n$ expansion, substituting $t$ for $\nu_1(t)$ in the Borel inversion
need not have an $o(1)$ effect. That second inversion is the source of
the cascade in Section~\ref{sec:resonance}.

\subsection{A family where the naive Lambert approximation fails}
Fix $\eta>0$ and take
\begin{equation}
 L(\ee^s)=\exp\left(\frac{s}{(\log s)^\eta}\right)
 \quad\hbox{for sufficiently large }s.
 \label{eq:fast-slow}
\end{equation}
The admissibility and second-derivative conditions hold eventually.
With $\ell=\log s$, direct differentiation gives
\[
 h=\frac{s}{\ell^\eta},\quad
 \log h=\ell-\eta\log\ell,\quad
 h'\log h=\ell^{1-\eta}+o(1).
\]
Therefore
\begin{equation}
 b(\ee^s)=\frac{s}{\ell^\eta}-\ell+\eta\log\ell
            -\ell^{1-\eta}+o(1).
 \label{eq:fast-slow-b}
\end{equation}
It follows that
\begin{equation}
 \frac{\ee^{b(n)}}{\ee^{W(L(n))}}
 =\exp\left(-\bigl(\log\log n\bigr)^{1-\eta}+o(1)\right).
 \label{eq:naive-failure}
\end{equation}
For $\eta>1$ the naive replacement happens to be valid at root precision.
At $\eta=1$ it loses the factor $\ee^{-1}$, and for $0<\eta<1$ its
relative error is unbounded. In particular, at $\eta=1$,
\begin{equation}
 u_n^{1/n}\sim
 \frac{a}{\ee}\,
 \exp\left(\frac{\log n}{\log\log n}\right)
 \frac{\log\log n}{\log n}.
 \label{eq:eta-one}
\end{equation}
This demonstrates why a theorem formulated only with $W(L(n))$ would
be false for the stated general class.

\section{A cascade of Borel-growth resonances}
\label{sec:resonance}

We now compute every term which survives at additive $o(1)$ precision
in the logarithm of the ordinary Borel normalizer for
\eqref{eq:stretch-root}. This is a finite transseries at each fixed
$\gamma$, but its number of terms increases without bound as
$\gamma\uparrow1$.

For the moment keep $\gamma\in(0,1)$ fixed, write
\begin{equation}
 S=\log(at),\qquad
 c_k(\gamma)=\frac1k\binom{k\gamma}{k-1},\qquad
 K_\gamma=\left\lfloor\frac1{1-\gamma}\right\rfloor.
 \label{eq:ck}
\end{equation}
Generalized binomial coefficients have their polynomial meaning here.
The amplitude $a=c_1$ is not related to the coefficient
$c_k(\gamma)$; the latter always carries its argument to avoid confusion.

\begin{theorem}[All nonvanishing logarithmic growth terms]
\label{thm:cascade}
For the stretched-logarithmic family and exponentially controlled
positive amplitudes,
\begin{align}
 \log\nu_1(t)
 &=S+\sum_{k=1}^{K_\gamma}
 c_k(\gamma)S^{1-k(1-\gamma)}-\gamma\log S+o(1),
 \label{eq:nu-cascade}\\
 \log\Bcal_1(t)
 &\sim\frac{at}{S^\gamma}
 \exp\left(\sum_{k=1}^{K_\gamma}
 c_k(\gamma)S^{1-k(1-\gamma)}\right).
 \label{eq:B-cascade}
\end{align}
The theorem permits all the slope perturbations in
\eqref{eq:linear-perturbation}.
\end{theorem}
\begin{proof}
Put $y=\log\nu_1(t)$. The defining equation for $\nu_1$ and
\eqref{eq:stretch-root} give
\begin{equation}
 y=S+y^\gamma-\gamma\log y+o(1).
 \label{eq:y-reverse}
\end{equation}
Let $Y$ be the unique large solution of $Y=S+Y^\gamma$.
With $w=Y/S$ and $T=S^{\gamma-1}$, its equation is
$w=1+Tw^\gamma$. The analytic implicit-function theorem applies near
$(T,w)=(0,1)$, and classical Lagrange inversion gives the convergent
expansion
\begin{equation}
 w=1+\sum_{k\ge1}\frac1k\binom{k\gamma}{k-1}T^k.
 \label{eq:scalar-reversion}
\end{equation}
Retaining terms through $K_\gamma$ leaves an error
$O(S^{1-(K_\gamma+1)(1-\gamma)})=o(1)$ in $Y$.

We also have
\begin{equation}
 y=Y-\gamma\log S+o(1).
 \label{eq:log-perturb}
\end{equation}
To check this at the required additive precision, substitute
$Y-\gamma\log S$ in $z-z^\gamma+\gamma\log z$.
Using $Y/S\to1$ and
$S^{\gamma-1}\log S\to0$, the residual from $S$ is $o(1)$.
The derivative of this function tends to one, so the residual produces
an $o(1)$ error in its large inverse. This proves
\eqref{eq:log-perturb}, and then \eqref{eq:nu-cascade}.
Exponentiate and use Theorem~\ref{thm:borel} to obtain
\eqref{eq:B-cascade}.
\end{proof}

\begin{corollary}[Resonant constants]
\label{cor:resonances}
At $\gamma=1-1/m$, $m\ge2$, the last term in
\eqref{eq:nu-cascade} is the constant $1/m$. Hence
\begin{equation}
 \log\Bcal_1(t)\sim
 \ee^{1/m}\frac{at}{S^{1-1/m}}
 \exp\left(\sum_{k=1}^{m-1}
 c_k(1-1/m)S^{1-k/m}\right).
 \label{eq:resonant-constant}
\end{equation}
\end{corollary}
\begin{proof}
At this parameter value $1-m(1-\gamma)=0$ and
$c_m(\gamma)=m^{-1}\binom{m-1}{m-1}=1/m$.
\end{proof}

For unit amplitudes, so that $S=\log t$, the first examples are
\begin{align}
 0<\gamma<\tfrac12:\quad
 &\log\Bcal_1(t)\sim tS^{-\gamma}\ee^{S^\gamma},
 \label{eq:first-region}\\
 \gamma=\tfrac12:\quad
 &\log\Bcal_1(t)\sim\ee^{1/2}tS^{-1/2}\ee^{\sqrt S},
 \label{eq:half}\\
 \gamma=\tfrac23:\quad
 &\log\Bcal_1(t)\sim\ee^{1/3}tS^{-2/3}
       \exp\left(S^{2/3}+\tfrac23S^{1/3}\right),
 \label{eq:two-thirds}\\
 \gamma=\tfrac34:\quad
 &\log\Bcal_1(t)\sim\ee^{1/4}tS^{-3/4}
       \exp\left(S^{3/4}+\tfrac34S^{1/2}+\tfrac{15}{32}S^{1/4}\right).
 \label{eq:three-quarters}
\end{align}
The term called a resonance here is a zero exponent in the logarithmic
normalizer. It is not a collision of Borel singularities: the Borel
transforms are entire. Omitting the constant term changes an asymptotic
equivalent into an incorrect one even though it does not change the
order or the infinite-type classification.

\subsection{Every factorial normalization}
The same argument, with
$R=\sigma^{-1}\log(at)$, gives the following more general formula:
\begin{equation}
 \log\Bcal_\sigma(t)\sim
 \sigma(at)^{1/\sigma}R^{-\gamma/\sigma}
 \exp\left(\sum_{k=1}^{K_\gamma}
 \frac{c_k(\gamma)}{\sigma^k}R^{1-k(1-\gamma)}\right).
 \label{eq:sigma-cascade}
\end{equation}
Indeed, $y=\log\nu_\sigma$ now obeys
$y=R+\sigma^{-1}y^\gamma-(\gamma/\sigma)\log y+o(1)$, and its
analytic core has $T=\sigma^{-1}R^{\gamma-1}$ in
\eqref{eq:scalar-reversion}. At $\gamma=1-1/m$, the corresponding
constant factor is $\exp(1/(m\sigma^m))$. This formula uses
$(n!)^\sigma$; for the $\Gamma(1+\sigma n)$ normalization one must
first make the argument change in \eqref{eq:gamma-growth}.

\subsection{Uniform transition through a resonance}
The constant-factor jump in a fixed-parameter expansion has a smooth
transition when the parameter is allowed to move. We state it for a
fully specified family to make uniformity unambiguous:
\begin{equation}
 \lambda_j(\gamma)=j\exp((\log(j+1))^\gamma),\qquad c_j=1,
 \quad 0<\gamma<1.
 \label{eq:uniform-family}
\end{equation}
Write $\Bcal_{1,\gamma}$ for its ordinary Borel transform.

\begin{lemma}[Uniformity on compact parameter intervals]
\label{lem:gamma-uniform}
For \eqref{eq:uniform-family}, the coefficient-root theorem and the
Borel-growth theorem hold uniformly for $\gamma$ in each compact
subinterval $J\subset(0,1)$. On $J$, the expansion
\begin{equation}
 b_\gamma(\ee^s)=s^\gamma-\gamma\log s+o(1)
 \label{eq:uniform-b}
\end{equation}
also holds uniformly.
\end{lemma}
\begin{proof}
Let $\gamma_-\le\gamma\le\gamma_+$ describe $J$.
The profile $L_\gamma(x)=\exp((\log(x+1))^\gamma)$ is admissible.
For convexity, set $\ell=\log(x+1)$; then
\[
 (xL_\gamma(x))''=
 \frac{\gamma L_\gamma(x)\ell^{\gamma-2}}{(x+1)^2}
 \left(\ell(x+2)+x(\gamma\ell^\gamma+\gamma-1)\right)>0.
\]
The last sign follows from
$\ell(x+2)-x>0$, using $\log(1+x)\ge x/(1+x)$, and the remaining
$\gamma$-terms are positive. Moreover
\[
 \sup_{\gamma\in J}\frac{xL_\gamma'(x)}{L_\gamma(x)}
 =O((\log x)^{\gamma_+-1})\longrightarrow0,
\]
and $\inf_{\gamma\in J}\log L_\gamma(x)\to\infty$.
Thus $b_\gamma\to\infty$, $b_\gamma=o(\log x)$, and all the compact
profile estimates in Lemma~\ref{lem:profile} are uniform. Its coercivity
proof is uniform as well: the same eventual bound on the elasticity
works for every $\gamma\in J$, and the bounded-index slopes are uniformly
bounded. A contrary sequence $\gamma_n\in J$ would therefore contradict
exactly that proof. This establishes the uniform root law.

In Lemma~\ref{lem:entire-transfer}, the coefficient errors
$\eta_n(\gamma)$ consequently tend to zero uniformly. The derivative
bound for $\log f_\gamma$ and the small, central, and large index
estimates are all uniform; the corresponding $N(t,\gamma)$ tend to
infinity uniformly. This proves uniform Borel transfer.
Finally, in Proposition~\ref{prop:moving},
$h_\gamma(s)=s^\gamma+o(1)$ with the same uniform derivative estimates;
$h_\gamma'\log h_\gamma=O(s^{\gamma_+-1}\log s)=o(1)$.
The Taylor remainder is uniform, proving \eqref{eq:uniform-b}.
\end{proof}

\begin{theorem}[A double-scaling resonance profile]
\label{thm:resonance-layer}
Fix $m\ge2$, put $S=\log t$, and let
\begin{equation}
 \gamma(t,\theta)=1-\frac1m+\frac{\theta}{\log S}.
 \label{eq:moving-gamma}
\end{equation}
For $\theta$ in every fixed compact real interval,
\begin{equation}
 \frac{\log\Bcal_{1,\gamma(t,\theta)}(t)}
 {tS^{-\gamma(t,\theta)}
  \exp\!\left(\sum_{k=1}^{m-1}c_k(\gamma(t,\theta))
         S^{1-k(1-\gamma(t,\theta))}\right)}
 \longrightarrow \exp\left(\frac{\ee^{m\theta}}m\right)
 \label{eq:resonance-layer}
\end{equation}
uniformly in $\theta$.
\end{theorem}
\begin{proof}
For large $t$, all the parameters lie in a fixed compact subinterval
of $(0,1)$, so Lemma~\ref{lem:gamma-uniform} applies. The analytic
implicit-function expansion \eqref{eq:scalar-reversion} has a uniform
radius and bounded remainders on a sufficiently small closed parameter
interval about $1-1/m$. Truncate it after the $m$th term. The error in
$Y$ is
\[
 O\!\left(S^{1-(m+1)(1-\gamma(t,\theta))}\right)
 =O\!\left(S^{-1/m}\ee^{(m+1)\theta}\right)=o(1)
\]
uniformly for bounded $\theta$. The logarithmic perturbation estimate
is uniform too. The $m$th retained term is
\[
 c_m(\gamma(t,\theta))S^{1-m(1-\gamma(t,\theta))}
 =c_m(\gamma(t,\theta))\ee^{m\theta}
 \longrightarrow\frac{\ee^{m\theta}}m.
\]
Subtract the first $m-1$ terms, exponentiate, and apply the uniform Borel
growth theorem. This proves \eqref{eq:resonance-layer}.
\end{proof}

At $\theta=0$ this gives the factor $\ee^{1/m}$. The leading terms in
the denominator must be evaluated at the \emph{moving} value of
$\gamma$; replacing them by their limiting parameter values changes
unbounded terms. The theorem asserts uniformity for bounded $\theta$,
not a simultaneous limit with $|\theta|\to\infty$ or $m\to\infty$.

\section{Exact checks and numerical diagnostics}
\label{sec:computations}

The supplied program separates exact finite verification from
floating-point asymptotic diagnostics. None of the latter is used as a
premise of a theorem.

\subsection{Independent rational coefficient computations}
For rational inputs, coefficients are generated without a symbolic
algebra system. If
$E_j(q)=\exp(\lambda_j\Uhat(q))=\sum_{r\ge0}e_{j,r}q^r$, then
\begin{equation}
 e_{j,0}=1,\qquad
 e_{j,m}=\frac{\lambda_j}{m}
            \sum_{r=1}^m r u_r e_{j,m-r},\qquad
 u_n=\sum_{j=1}^n c_j e_{j,n-j}.
 \label{eq:code-recurrence}
\end{equation}
This follows by differentiating $E_j$, and supplies a triangular
recurrence independent of enumerating Lagrange compositions.

The recorded run uses $\lambda_j=jH_j$. It compares the recurrence
with direct composition enumeration, checks convex composition bounds,
checks the finite lower and upper inequalities using exact fractions,
and repeats a coefficient comparison for the nonconstant amplitudes
$c_j=(j+1)/2$. For the latter,
$1\le c_j\le2^{j-1}$, so \eqref{eq:amplitudes} holds with $a=1$ and
$D=\log2$. Separately, the coefficients of
$w=1+Tw^\gamma$ are checked through degree nine for
$\gamma=1-1/m$, $2\le m\le8$, including $c_m(\gamma)=1/m$.

The recorded output reports \textbf{467 exact assertions passed}, with
rational feedback coefficients computed through order 24. These are
finite arithmetic checks, not a formal proof of the limit theorems.

\subsection{Finite coefficient envelopes}
The first table uses unit amplitudes and the harmonic slopes
\eqref{eq:harmonic-slopes}. Its lower and upper columns are the logarithms
of the two sides of \eqref{eq:finite-u-bound}, divided by $n$.
The scale $b_n$ is computed with the smooth harmonic interpolation, and
$m_{\rm low}$ maximizes the lower envelope.
\begin{center}
\small
\begin{tabular}{rrrrr}
\toprule
$n$ & lower $\log u_n/n$ & $b_n$ & upper $\log u_n/n$ & scaled maximizer\\
\midrule
100 & 1.383463 & 1.252376 & 1.999557 & 0.833379\\
1,000 & 1.618715 & 1.486732 & 2.270770 & 0.855436\\
10,000 & 1.783647 & 1.665473 & 2.433987 & 0.874808\\
100,000 & 1.915549 & 1.809832 & 2.559553 & 0.889537\\
\bottomrule
\end{tabular}

\end{center}
The two envelopes need not lie on opposite sides of the asymptotic
normalizer $b_n$: they enclose $\log u_n/n$, not its limit expression.
In fact the lower column already exceeds $b_n$ here. The convergence
is very slow, and the upper estimate deliberately counts all compositions
using a common maximal slope sum. At the displayed sizes the table does
not identify the limiting constant by itself. These floating-point values
are not outward-rounded interval certificates.

\subsection{Resonant normalizer diagnostics}
For the next table use the large-$x$ profile
$L(x)=\exp((\log x)^\gamma)$ and let $S=\log t$.
Instead of forming enormous numbers $t$ or $\nu_1(t)$, solve directly for
$b=\log(\nu_1(t)/t)$:
\begin{equation}
 b+\log b=\bigl(S+b-\log(1+b)\bigr)^\gamma.
 \label{eq:numerical-b}
\end{equation}
At $\gamma=1-1/m$, the displayed quantity is
\[
 \exp\left(b+\gamma\log S-
        \sum_{k=1}^{m-1}c_k(\gamma)S^{1-k/m}\right),
\]
whose limit is $\ee^{1/m}$. The calculation uses 75 decimal working
digits.
\begin{center}
\small
\begin{tabular}{rrrr}
\toprule
$\gamma$ & $\log t$ & normalizer ratio & limit $e^{1/m}$\\
\midrule
$1/2$ & $10^4$ & 1.6446049088 & 1.6487212707\\
$1/2$ & $10^8$ & 1.6486602121 & 1.6487212707\\
$1/2$ & $10^12$ & 1.6487206526 & 1.6487212707\\
$2/3$ & $10^4$ & 0.9366736042 & 1.3956124251\\
$2/3$ & $10^8$ & 1.3456993528 & 1.3956124251\\
$2/3$ & $10^12$ & 1.3921133863 & 1.3956124251\\
$3/4$ & $10^4$ & 0.4031221486 & 1.2840254167\\
$3/4$ & $10^8$ & 1.0358703247 & 1.2840254167\\
$3/4$ & $10^12$ & 1.2439118048 & 1.2840254167\\
\bottomrule
\end{tabular}

\end{center}
These are diagnostics of the \emph{implicit growth normalizer}, not
numerical evaluations of $\Bcal_1(t)$ or proofs that the same finite
values approximate it closely. The JSON output also contains six
moving-parameter normalizer checks for Theorem~\ref{thm:resonance-layer},
computed at 100 decimal working digits; these have the same limitation. Even $S=10^{12}$ leaves a visible
subleading correction for $\gamma=3/4$. Retaining the negative-power
terms in \eqref{eq:scalar-reversion} and the next moving-argument
correction would substantially improve numerical convergence.

\subsection{Reproduction}
From the archive root, the commands are
\begin{quote}\ttfamily\small
python -m pip install -r requirements.txt\\
python code/verify.py\\
pdflatex -interaction=nonstopmode -halt-on-error near\_linear\_feedback.tex\\
pdflatex -interaction=nonstopmode -halt-on-error near\_linear\_feedback.tex
\end{quote}
Exact arithmetic uses the standard Python library; NumPy and SciPy
supply the large finite envelopes, and mpmath supplies high-precision
normalizer calculations. The script writes both machine-readable output
and the two tables included above. A third LaTeX pass may be used when
section or table placement has changed.

\section{Ten further research questions}
\label{sec:future}

\subsection{A multiplicative coefficient equivalent}
Determine the $\exp(o(n))$ factor in
$u_n=(a\ee^{b_n})^n\exp(o(n))$, first for $\lambda_j=jH_j$ and unit
amplitudes. The finite upper envelope ignores the exact distribution of
slope sums across compositions. A solution will require finer entropy
estimates and a saddle correction, not just a higher-order expansion of
$b_n$. A useful intermediate target is the first sublinear term of
$\log u_n$ with a proved error smaller than that term.

\subsection{A local limit theorem for composition length}
The rate function has $I''(1)=1$, suggesting fluctuations of order
$x_n/\sqrt n$ around a suitably corrected value of $n-K_n+1$.
That scale and a Gaussian law are conjectural here. Terms discarded as
$o(n)$ can shift the centre by much more than the fluctuation width.
Determine the correct centre and variance before asserting a central
limit theorem, and then seek a uniform lattice local limit.

\subsection{Condensation of the primitive actions}
Does a typical positively weighted Lagrange composition have one
exceptional part containing almost all its excess? The present length
law neither proves nor disproves this. For $L(x)=\exp((\log x)^\gamma)$,
splitting an action changes $L$ by a relative amount of order
$(\log x)^{\gamma-1}$ on fixed-ratio scales. This points to a possible
sub-$n$ rate for condensation. Establish that rate, identify competing
entropy, and distinguish unique condensation from several macroscopic
parts.

\subsection{Minimal hypotheses and oscillatory slow variation}
Remove the differentiability and convexity assumptions while retaining a
correct normalizer. Convexity provides a sharp exceptional-part upper
bound; without it, the relevant object may be an upper convex envelope
of the discrete slopes rather than $jL(j)$ itself. Seek necessary and
sufficient conditions under which that envelope preserves coefficient-
root growth, and analyze slowly varying factors with long plateaux or
oscillatory local elasticity.

\subsection{Sparse actions and much smaller amplitudes}
Allow $c_j=0$, noninteger actions, or amplitudes with
$\log c_j$ more negative than order $j$. The proof here uses the ability
to fill almost every remaining unit of degree with action-one factors,
and it needs at least exponentially sized exceptional-action amplitudes.
Determine how support gaps and amplitude suppression alter the optimal
length, the Lambert balance, and the rate function. A first concrete
case is $c_j=\exp(-j\log\log(j+2))$.

\subsection{Signed feedback and vector systems}
For signed or complex amplitudes the positive coefficient model is only
an absolute majorant. Develop verifiable noncancellation conditions that
recover lower growth bounds. For vector equations, determine whether a
one-dimensional effective slope emerges from positive matrix feedback,
or whether several interacting large actions produce a different
variational problem. Any extension must state its infinite-tail
conditions rather than invoke a finite-dimensional theorem informally.

\subsection{Growth away from the positive Borel ray}
Determine the directional growth of $\Bcal_\sigma$ on nonpositive rays
and in sectors. Positivity fixes maximum modulus but does not prevent
cancellation on other rays. Compare any resulting Laplace sums with the
left-sector analytic fixed points constructed in the repository feedback
article. Neither an entire Borel transform nor coefficient Gevrey bounds
alone proves that the two analytic constructions coincide.

\subsection{Compatible acceleration and optimal truncation}
Use the scale $\nu_1$ to construct an acceleration, if possible, which
respects multiplication and the nonlinear feedback equation after an
explicit analytic interpretation of its countable kernel. The damping
threshold in Proposition~\ref{prop:damping} is only a necessary guide.
For a specified sectorial realization, also establish uniform remainder
bounds and an actual optimal-truncation theorem. Coefficient-root growth
alone does not certify errors of truncated analytic solutions.

\subsection{Inversion and parameter limits beyond fixed resonances}
Determine the exact root-growth scale of the compositional inverse of
$\Uhat$. Equality of the least Gevrey index is too weak to decide whether
$a\ee^{b_n}$ survives reversion. Separately, extend
Theorem~\ref{thm:resonance-layer} to an increasing resonance index $m$,
or to the transition $\gamma\uparrow1$ at a rate depending on $t$.
That limit requires uniform control of a growing number of reversion
terms and should connect to the genuinely superlinear regime.

\subsection{Formal proofs and certified asymptotic computation}
A practical Lean sequence starts with the finite coefficient formula,
the discrete convex composition bound, and the exact amplitude
inequalities. The next stages are the compact profile lemma, its tail
coercivity, and the maximum-term transfer. These are substantial but
separable targets; no full transseries datatype is needed for the first
stage. On the numerical side, implement directed interval versions of
\eqref{eq:finite-u-bound}, certified roots for \eqref{eq:b-definition},
and explicit moduli for the $o(1)$ terms. Keep finite interval enclosures,
asymptotic theorems, and Lean certificates as distinct output types.

\section{Conclusions}

The near-linear regime is not a degenerate version of the superlinear
Gevrey calculation. It has its own self-consistent Lambert scale:
\[
 b_n\ee^{b_n}=L\left(\frac n{1+b_n}\right),\qquad
 u_n^{1/n}\sim a\ee^{b_n}.
\]
The coefficient theorem tolerates substantial amplitude variability and
bounded linear slope errors. At the same precision it provides a universal
large-deviation law for Lagrange composition length, with a rate function
independent of the slowly varying profile.

The scale transfers to entire Borel growth through one more nonlinear
inversion. That transfer distinguishes zero Gevrey type at every positive
order from positive-direction summability: all the normalized transforms
are entire but too large for their standard Laplace kernels. For
stretched-logarithmic amplification the inversion yields a finite,
explicit power--logarithmic expansion at each parameter, an infinite
sequence of resonant constants, and a uniform transition profile through
each fixed resonance. A faster slowly varying example shows why the
moving argument in the initial Lambert equation is indispensable.

These conclusions answer the inspected repository's near-linear question
at coefficient-root and transform-growth precision under the stated
hypotheses. A full multiplicative coefficient equivalent, general signed
or sparse systems, and a compatible acceleration theory remain separate
research tasks rather than consequences asserted without proof.

\appendix
\section{Source snapshot and claim boundaries}
\label{app:sources}

\subsection{Pinned repository source}
The initial survey used the repository root tree
\begin{center}\small\ttfamily
 a795fcffa4b3ece22761d81fbed3c43be8b81795.
\end{center}
The unchanged predecessor was rechecked at commit
\begin{center}\small\ttfamily
04473354a0f3edff2d6c365caf26170ca6b88735.
\end{center}
The initial transseries project tree was
\begin{center}\small\ttfamily
 e37727630fd519201ceb2b2f24fdc71b238873d9.
\end{center}
The direct predecessor is the file \texttt{article.tex} in
\begin{quote}\small\ttfamily\raggedright
Analysis/Transseries/docs/series-and-transseries/\\
Exponential\_Feedback\_Regularity\_Classification/
\end{quote}
Its blob hash is
\begin{center}\small\ttfamily
 c6a83d71973a12fe7ed970d2abc56e2b49e6808f.
\end{center}
The contiguous source range 1400--1605 contains the further-research
section, including the near-linear, amplitude, and summation questions.
The source range 350--530 contains the positive coefficient formula,
its proof, and the convergence criterion. The group README initially
recorded six unmerged arrivals of 29 September 2026 and recorded twelve
at the final recheck, after a second intake. Both versions describe the
arrivals as lacking claim-by-claim review and Lean formalization.
We therefore reprove the needed finite algebra rather than treating the
predecessor as a formally verified dependency.

\subsection{Scope of the repository comparison}
The review used the current transseries project README, the group README
and its descriptions of the initial six arrivals, the subsequent intake
descriptions at the final recheck, and the relevant contiguous parts of
the feedback article. It was not a line-by-line audit of the
entire canonical and companion volumes or every repository file.
No repository branch was changed. The proposed placement of this package
is a new research directory beside the predecessor, pending mathematical
review and normal repository intake.

\subsection{Status matrix}
\begin{center}\small
\begin{tabular}{@{}p{0.38\linewidth}p{0.55\linewidth}@{}}
\toprule
Item & Status and boundary\\
\midrule
Formal coefficient formula & Classical tool; reproved with amplitudes and checked at finite orders.\\[3pt]
Coefficient-root equivalent & Proved under admissibility, positivity, and exponential amplitude bounds; not a full coefficient equivalent.\\[3pt]
Length large deviations & Proved with speed $n$; does not identify a largest-action or Gaussian law.\\[3pt]
Entire Borel growth & Proved on the positive ray and hence for maximum modulus; not a directional classification.\\[3pt]
Positive summation obstruction & Proved for every standard finite Gevrey order; not an obstruction to all generalized summation.\\[3pt]
Resonance cascade and layer & Proved for the stated families; fixed $m$, compact moving parameters.\\[3pt]
Computation & 467 exact finite assertions plus nonrigorous floating-point diagnostics; no new Lean proof.\\[3pt]
Novelty & Extension of an explicit inspected repository question; external publication priority not exhaustively established.\\
\bottomrule
\end{tabular}
\end{center}

\section{A proof-audit guide}
\label{app:audit}

The central technical point is Lemma~\ref{lem:profile}. Its compact
identity is exact before taking limits, and its coercivity argument is
needed to discard the exponentially many but poorly weighted
compositions. Counting all compositions by $2^n$ before localization
would destroy the constant at coefficient-root precision. The amplitude
hypothesis is useful precisely because its logarithmic cost is at most
$D(n-K_n)$, which is $o(n)$ in the localized regime.

The next sensitive step is the transfer to Borel growth. An error
$\exp(o(n))$ in coefficients is compatible with an equivalent for
$\log\Bcal_\sigma(t)$, but not with an equivalent for
$\Bcal_\sigma(t)$ itself. All the displays in
Section~\ref{sec:resonance} deliberately concern the logarithm of the
Borel transform. Exponentiating those displays a second time without a
stronger error estimate would be invalid.

Finally, when a reversion term has exponent zero it must be retained
before exponentiating the logarithmic normalizer. This is the source of
the constants in Corollary~\ref{cor:resonances}, and it is why the index
$K_\gamma$ includes the equality case. The moving-parameter theorem
requires uniform estimates from the coefficient theorem onward; its
proof does not merely insert a moving parameter into a pointwise
asymptotic statement.

\clearpage
\phantomsection
\addcontentsline{toc}{section}{References}
\begin{thebibliography}{9}

\bibitem{repo-feedback}
Vladimir Reshetnikov, ProveIt research repository.
\emph{A Sharp Regularity Classification for Countable
Exponential-Feedback Transseries}. Research article filed 29 September
2026; source snapshot and precise path in Appendix~\ref{app:sources}.
\href{https://github.com/VladimirReshetnikov/ProveIt/blob/04473354a0f3edff2d6c365caf26170ca6b88735/Analysis/Transseries/docs/series-and-transseries/Exponential_Feedback_Regularity_Classification/article.tex}{Pinned TeX source}.
The repository owner is given for attribution of the source location;
no independent authorship or peer-review claim is intended.

\bibitem{gessel}
Ira M. Gessel.
Lagrange inversion.
\emph{Journal of Combinatorial Theory, Series A} \textbf{144}
(2016), 212--249.
\href{https://doi.org/10.1016/j.jcta.2016.06.018}{doi:10.1016/j.jcta.2016.06.018};
\href{https://arxiv.org/abs/1609.05988}{arXiv:1609.05988}.

\bibitem{jkt}
Sabine Jansen, Tobias Kuna, and Dimitrios Tsagkarogiannis.
Lagrange inversion and combinatorial species with uncountable color
palette.
\emph{Annales Henri Poincar\'e} (2021).
\href{https://doi.org/10.1007/s00023-020-01013-0}{doi:10.1007/s00023-020-01013-0};
\href{https://arxiv.org/abs/2008.10862}{arXiv:2008.10862}.

\bibitem{bgt}
N. H. Bingham, C. M. Goldie, and J. L. Teugels.
\emph{Regular Variation}.
Encyclopedia of Mathematics and its Applications, vol.~27.
Cambridge University Press, 1987.
Background reference; the differentiable estimates needed here are
proved in the text.

\bibitem{sauzin}
David Sauzin.
\emph{Introduction to 1-summability and resurgence}.
2014.
\href{https://arxiv.org/abs/1405.0356}{arXiv:1405.0356}.

\bibitem{nikolaev}
Nikita Nikolaev.
\emph{Gevrey Asymptotic Implicit Function Theorem}.
2021, version 2.
\href{https://arxiv.org/abs/2112.08792}{arXiv:2112.08792}.

\bibitem{dlmf}
NIST Digital Library of Mathematical Functions.
\S4.13, Lambert $W$-function.
\url{https://dlmf.nist.gov/4.13}.
Accessed 29 September 2026.

\end{thebibliography}
\end{document}
